% FILE: 04_implementation_surface.tex
% PROJECT: gaps — Structural Gap Registry and Closure Tracking

\section{Implementation Surface}
\label{sec:gaps-implementation}

\subsection{Source Files}
\label{sec:gaps-implementation-source}

The implementation surface of the gap registry consists entirely of authored Markdown
documents. No Rust source files, ontology files, receipt files, or test fixtures reside
in the \texttt{gaps/} directory itself. The directory is a documentary layer; it
describes and analyzes implementation surfaces located elsewhere in the
\texttt{process-intelligence} foundry and in the adjacent \texttt{wasm4pm} repository.

The five source files are:

\begin{description}
  \item[\texttt{GAP\_REGISTER.md}] The master index. Provides the canonical status table
        (eight gaps, severity, status, and priority), a two-line problem statement for
        each gap, and a resolution description. The register is the authority surface;
        any discrepancy between the register and an individual gap document should be
        treated as a discrepancy requiring a resolution addendum, not silently resolved
        in favor of either source.

  \item[\texttt{GAP\_001\_COMPAT\_WASM\_BRIDGE.md}] The most detailed individual gap
        document. Contains the full Cargo.toml evidence block, three Rust function
        signatures from wasm4pm algorithm files (\texttt{ilp\_discovery.rs},
        \texttt{process\_tree.rs}, \texttt{performance\_dfg.rs}), the complete wasm4pm
        error enum from \texttt{wasm4pm-types/src/error.rs}, the two-universe ASCII
        diagram, the five-condition remediation checklist, and the authorization clause
        blocking downstream refactors until the research program approves the remediation
        path.

  \item[\texttt{GAP\_002\_OR\_JOIN\_AMBIGUITY.md}] Documents the BPMN OR-Join
        non-determinism problem and the v30.1.1 Smart-Completion resolution. Includes
        the formal mathematical policy expressed in LaTeX notation. Shorter than GAP\_001
        because the resolution is documented rather than pending.

  \item[\texttt{GAP\_PRIORITY\_MATRIX.md}] The dependency and sequencing analysis.
        Provides the formal dependency graph as an ASCII directed acyclic graph, the
        priority matrix table, the P1/P2/P3 closure sequences, the critical path
        (\texttt{GAP\_007 $\rightarrow$ GAP\_001 $\rightarrow$ GAP\_002}), and the
        blocking-claims table that maps each board-level or research-level claim to the
        gap that must be closed before the claim can be made.

  \item[\texttt{GAP\_CLOSURE\_WORKFLOW.md}] The governance workflow document under the
        V30.1.1 Epoch. Provides the three constitutional directives, a three-phase
        execution plan (Phase 1: P1 gaps; Phase 2: P2 and P3 gaps; with specific
        remediation steps for each gap), and a description of the SWARM Verification
        Guard architecture. The SWARM Verification Guard (Drift Sentry, Telemetry
        Auditor, Ledger Custodian, Alignment Referee, Stream Director) is described as
        a five-agent monitoring system. No corroborating evidence of its implementation
        exists in the \texttt{gaps/} directory; this description is noted as
        potentially aspirational architecture. \textbf{INTERPRETATION:} The research
        program treats this section as a design intent statement, not an implemented
        system claim.
\end{description}

\subsection{Commands}
\label{sec:gaps-implementation-commands}

No CLI commands are defined within the gap registry itself. The registry is a read-write
documentary surface maintained by human or agent authorship following the governance
rules of the V30.1.1 Epoch. Gap records are authored in Markdown following the
template structure visible in the five existing documents: gap ID, severity, status,
discovered date, source reference, summary, expected capability, observed capability
(with code citations), impact assessment, and authorization clause.

The commands that would constitute gap closure actions are located in the remediating
projects: Cargo commands in wasm4pm for dependency addition, Rust source edits in
wasm4pm for algorithm signature changes, and fixture replacements in wasm4pm-compat for
GAP\_008. The gap registry does not execute these commands; it authorizes and tracks them.

\subsection{Data and Ontology Surfaces}
\label{sec:gaps-implementation-data}

The gap registry does not define ontology files, TTL graphs, or structured data schemas.
Its data surface is the gap record format: a structured Markdown template with mandatory
fields (gap ID, severity, status, discovered date) and recommended sections
(summary, expected capability, observed capability, impact assessment, remediation path,
authorization). This format is consistent across all five existing documents and
constitutes the implicit schema of the registry.

The dependency graph in \texttt{GAP\_PRIORITY\_MATRIX.md} is a data structure in
documentary form — an ASCII directed acyclic graph encoding the partial order over gap
closures. The blocking-claims table is a relational structure mapping claim strings to
gap identifiers. Neither is machine-readable in its current form; both are human-auditable
documentation intended to govern the sequencing of remediation work.

\subsection{Templates and Rendered Artifacts}
\label{sec:gaps-implementation-templates}

The gap record template implicit in the existing documents provides the pattern for
future gap additions. Each new gap document must include: a formal gap ID in the series
GAP\_001 through GAP\_N; a severity classification (CRITICAL, MAJOR, or MINOR); a
status (Open or Resolved with resolution reference); a discovery date and source
citation; a summary of the defect; an expected capability statement grounded in
existing doctrine; an observed capability statement with direct code citations; an
impact assessment table; a remediation path with specific conditions; and an
authorization clause.

Resolution addendums follow the same template with a header marking them as addendums
and a date. Gaps are never edited to remove their Open status; the status field is
updated in the master \texttt{GAP\_REGISTER.md} and the resolution addendum is appended
to the individual gap document. This is the immutability doctrine applied to documentary
artifacts.
